% 第5章（本文を親の main.tex から \input してください）
% 画像は figures/chapter5/ に配置してください。

\chapter{一次史料と二次情報の意味的接続}

\section{はじめに：背景と対象}

4章では、HIMIKOモデルを中心に、史料の内容情報の構造化について論じられた。その中では、EntityInContextの概念など、史料中で言及されるエンティティと外部知識との接続についても扱った。ただし、そこでいう外部知識とは、あくまで史料において言及されたエンティティについてのコンテクストを補完するものであった。他方、歴史研究においては、史料の側の記述において参照されるわけではないものの、その内容を理解し解釈するために必要な関連文献や資料などの二次情報を常に参照しうる。むしろ、そうした二次情報を広範に参照し既存の学説や解釈の蓄積を踏まえた上でなければ、一次史料そのものを適切に読解し解釈することはできない。

すなわち、一次史料とそれに関する研究論文・書籍・資料等からなる広範な学術情報ネットワークが構築され、それを活用しながら解釈を行うことのできる環境が、歴史研究には不可欠である。もっとも、こうした学術情報ネットワークの構築に関しては従来、学術情報流通の文脈において多くの文献探索データベースや電子ジャーナル、デジタルアーカイヴの整備が進められてきたものの、実際に研究者が読解・解釈の対象とする狙いの史料、さらにはその特定の箇所についてどのような関連情報が存在するのかについては、研究者個人が持つ知識に基づいて、経験的に探索するほかなかった。

しばしば経験的な職人芸に依存する一次史料-二次情報接続の一つの例が、古注（スコリア）である。古注とは、主に古典文献の原文に対する近現代以前になされた注釈を指す。たとえば西洋古典であれば、アリストテレスの著作に対して、古代末期や中世になされた注釈がしられる。この古注が作成された時代には、現代のように精密なテクスト参照規則が定まっていなかったことから、特定の注が原文のどの箇所に該当するのかも、現代の正書法からみて必ずしも定かではない場合が少なくない。また、ある特定のテクストに対しどのような注釈者・注釈書が存在するのかということ自体、必ずしも情報が整備されているとは限らず、専門外の者には発見が困難である。

以上のように、まず注釈の存在を認識し、さらに原文と注の対応をつけるという作業は、専門的な訓練によって涵養される経験的なスキルなくしては困難である。もっとも、こうしたスキルそのものは人文学研究において重要なものであることは間違いない。一方で、一次史料の特定のテクストと、その注釈書、さらには注釈書中の特定の注の関係性を、一貫したデータとして記述し展開することができれば、情報の発見可能性と利用可能性は大幅に向上するはずである。

もっとも、こうした一次史料-二次情報接続による学術情報の発見・利用可能性の向上は、古注に限られるものではない。歴史研究においては、先行研究における史料解釈が参照されることが他の分野以上に重要になるが、現状では、どの文献においてどの史料が言及されているか、ということすら基礎的な情報でも十分に整備されておらず、関連ワードやキーワードによる網羅性の点では課題の大きい検索手法に依存せざるを得ない。

こうした状況を踏まえるに、一次史料と、注釈・研究文献・関連資料などを含む二次情報の接続もまた、デジタルヒストリーが取り組むべき重要な課題の一つであることは明らかであろう。また、一次史料がどのような関連情報のネットワークの中で読み、解釈しているのか、を記録・展開するという観点からも、本研究の根幹的なテーマである歴史研究における解釈とプロセスの構造化にとっても、欠かせない要素の一つである。

そこで本章では、こうした一次史料と二次情報の接続と構造化を実現するための概念的枠組みとデータモデルを検討し、データ構築に向けたワークシステムの設計について論じる。ただし、こうした史料と二次情報の接続について、いきなり歴史研究一般を対象とすることは困難であり、あくまで、対象を絞って論じることが必要になる。たとえば、上述の古注の例と、現代における注釈や研究文献の例では、実際の情報記述のあり方は大きく異なるはずである。また、対象となる一次史料の類型によっても、その解釈において参照される二次情報の扱われ方は大きく異なる。

そのため本章においては、検討の対象をあくまで同様にラテン語碑文とし、その研究プロセスにおいて参照される二次情報に絞る。碑文研究においては、古典テクスト研究におけるような古注はほとんど存在せず、研究者が参照するのは主に近現代のテクスト校訂や史料集、研究文献、考古学調査報告などである。

\section{問題設定：一次史料のアノテーションと文献からの史料言及抽出}

一次史料と二次情報の接続と一口にいっても、その研究プロセスの中で現れるその具体的な相貌は多様である。ただ実際にデータ構造化を進めるにあたっては、ある程度限定的な作業課題を設定する必要がある。その際、大きな方向性としては、（1）既存の文献や書籍といった二次情報から、一次史料参照情報を抽出する「二次情報→一次史料」という方向性、（2）狙いの一次史料にさらに言及したり、あるいは関連づけてトピックや概念を扱う二次情報を接続する「一次史料→二次情報」という方向性の二つが考えられるだろう。

これらの二つの方向性は、いずれも最終的には一次史料と二次情報の接続に連なるものではありつつ、具体的な情報処理タスクとしてみた場合、その内実は大きく異なるだろう。

まず、「二次情報→一次史料」においては、既存の文献・書籍というテクストデータから、「史料」というエンティティへの言及を抽出する、固有エンティティ抽出が中心的な課題となる。史料表記の正規化や名寄せといった問題に取り組む必要がある。また、史料言及の位置情報やコンテクストを保持することを考えれば、文献・書籍を章や節といった単位で細分化して構造化し参照可能にする必要が生じる。他方、「一次史料→二次情報」の場合には、一次史料に対するアノテーション手法の検討が主たる課題になるだろう。アノテーションとは、テクストや画像といったデータ（あるいはその特定箇所）について、追加の情報（メタデータ）を付加する作業を指す。一次史料に対するアノテーションは、人名や地域についてのエンティティ・リンキングを中心にデジタルヒューマニティーズの分野でも盛んに行われているが、これを書籍・文献等にさらに拡張することが主要な課題となる。

このように、データ構築プロセスの方向性によって、作業課題には大きな相違が生じる。そのため本章では以下それぞれについて個別に事例研究を行う。

一方で、両者に共通の課題も存在する。それは、一次史料と二次情報の「意味的接続」の実現である。一次史料と二次情報の関係を表現する最も単純な方法は、両者の間に「参照される」というリンクを設定することである。しかしながら、この二項関係は、接続の存在を示す一方で、参照が果たす機能を捉えるにはあまりに粗い。たとえば、同じ碑文番号への言及であっても、ある箇所では碑文の存在を示すだけであり、別の箇所ではラテン語本文を翻刻し、さらに別の箇所では年代を推定し歴史的議論の証拠として利用される。これらを同一の参照関係として扱えば、史料が学術的言説のなかで担う役割の差異は失われる。「意味的接続」は、一次史料と二次情報を単にURIで結ぶのではなく、どの箇所においてどの史料がどのような機能のもとで扱われているかを記述し、上述のような情報の喪失を防ぎ、参照関係の精緻な検索・比較・再利用可能性を担保するものである。

\section{二次研究文献からの史料言及抽出}

\subsection{既存の引用・参照モデルと残された課題}

\subsubsection{引用索引から意味ある参照関係へ}

従来の引用索引や計量書誌学では、研究成果相互の引用関係が主要な対象となってきた。この枠組みは研究動向や影響関係の分析には有効であるが、歴史学において一次史料がどのように用いられたかを直接には表現しない。二次文献から一次史料への参照は、先行研究への賛否や継承とは異なり、史料の提示や年代の推定、批判、解釈、証拠化といった複数の機能を含みうる。そのため、Publication cites Publication という一般的な引用ネットワークをそのまま Publication refers to Primary Source へ置き換えるだけでは不十分である。

\subsubsection{CIDOC CRM}

文化遺産情報の統合を目的とするCIDOC CRMは、物理対象・情報オブジェクト・文書のあいだの関係を記述するための基礎的な語彙を提供する。たとえばP67 refers to、P70 documents、P62 depictsなどを利用すれば、文書がある対象を参照・記録・描写することを表現できる。この点で、一次史料と研究文献を同一の知識グラフに配置する基礎として有用である。

しかしながら、これらのプロパティは広範な相互運用性を優先するため、狙いの学術的参照が何を行っているかを細かく区別しない。碑文本文の翻刻、年代の提示、異読の提案、歴史的議論の証拠化は、いずれも広い意味では「参照」または「記録」に含まれうるが、研究者が必要とする検索条件としては区別されるべきである。したがって、本章の課題はCIDOC CRMを排除することではなく、その上位の汎用的関係に対して、参照行為の領域固有の意味論を補うことにある。

\subsubsection{CiTO}

CiTO（Citation Typing Ontology）は、引用を単なる有向辺としてではなく、support、criticize、confirmなどの意味づけ関係として記述するためのオントロジーである。その発想は、本章が目指す意味的接続に直接的な示唆を与える。すなわち、リンクの存在に加えて、そのリンクが学術的言説のなかで果たす機能を明示する必要がある。

一方、CiTOの主要な想定は学術文献間の引用関係であり、歴史研究に特有の二次文献から一次史料への関係を網羅するものではない。碑文の翻刻、翻訳、複数間版間の番号の同定、欠損部の補塡、論争的出版状況、考古学的に検討するといった行為は、一般的な引用の修辞的機能だけでは捉えにくい。そこでCiTOの「引用関係を型付けする」という発想を継承しつつ、一次史料参照に適した関係タイプを別途設計する必要がある。

\subsubsection{HuCIと文脈保存}

HuCI（Humanities Citation Index）の提案は、人文学の引用が持つ曖昧性と文脈依存性を強調する。引用が支持的か批判的か、単なる言及か証拠としての利用か、記述的か分析的か、は引用元と引用先の書誌情報だけからは判定できない。判断には、引用が置かれた文章・節・注や前後の議論を参照する必要がある。HuCIが掲げる「Rich in Context」という観点は、一次史料と二次文献の意味的接続にも不可欠である。

この観点から導かれる第一の要件は、文献全体と史料全体を直接結ぶだけではなく、文献中の具体的な言及を独立したノードとして保持することである。第二の要件は、その言及を段落・脚注・節などの文書構造に位置づけ、必要に応じて周辺テクストへ戻れるようにすることである。第三の要件は、意味分類の結果だけでなく、その判断を支える根拠表現を保存することである。本章のモデルは、この三要件をMention、Relation Assertion、Textual Contextの結合によって実装する。

\subsection{一次史料参照の概念モデル}

\subsubsection{文献単位からMention単位へ}

本章のモデルにおける中心的な単位はMentionである。Mentionとは、二次文献の特定箇所に現れる、特定の一次史料への一回の言及を指す。Publicationから Primary Sourceへ直接プロパティを設定するモデルでは、同じ文献内の複数箇所に現れる参照が区別できず、本文と脚注の差異や、異なる参照機能、複数の議論文脈が失われてしまう。これに対してMentionを介在させれば、各言及に固有の属性と関係を付与できる。

各Mentionは、原文に現れる参照表記（raw reference）、正規化された表記（normalized reference）、本文または脚注という位置種別、開始・終了段落、注番号などを保持する。さらに、Mentionは参照対象となるInscription、参照表記を構成するCorpus Reference、文書構造上のParagraphまたはFootnote、そして意味分類を担うRelation Assertionへ接続される。この構造により、文献・言及箇所・参照表記・史料実体・意味タイプが相互に分離しつつ関連づけることができる。モデル上、MentionはWeb Annotation語彙の \texttt{oa:Annotation} のサブクラスとして表現し、周辺テクストは注釈本体（\texttt{oa:TextualBody}）として保持する。

\subsubsection{Relation Assertionと多値分類}

一次史料への一つの言及は、単一の機能だけを持つとは限らない。脚注にラテン語本文と訳文を併記し、その前後で年代と出土地を説明する場合、その言及は、翻刻文提供、翻訳提供、年代提示、出土地・所在記述を同時に満たしうる。したがって、参照タイプを相互排他的なクラスとして扱うことは適切ではない。本章では、一つのMentionに複数のRelation Assertionを関連づける多値分類を採用する。

Relation Assertionを独立したノードとする利点は、分類語彙だけでなく、判断根拠、生成主体、使用モデルやプロンプト、実行時刻、確信度、専門家による検証状態などを、将来的に付与できる点にもある。とくにLLMを利用する場合、分類結果を不可視の属性値として保存するのではなく、再検証可能な解釈行為の産物として扱う必要がある。これは、歴史知識が確実的事実だけでなく、史料と文脈に基づく解釈の履歴として記述するという本論文全体の方針とも整合する。

\subsubsection{文書構造と参照文脈}

参照の意味は、しばしば史料番号を含む一文だけでは確定しない。本文の段落が歴史的解釈を提示し、脚注が翻刻と異読番号を与える場合、両者は一つの議論単位を形成する。そこでモデルは、Source Document、Book Section、Paragraph、Footnoteなどの階層を保持し、Mentionを該当箇所へ接続する。これにより、検索結果から原文の深いリンクへ遷移し、意味分類の妥当性を人間が確認できる。

文脈の範囲は固定的ではなく、参照表記の周辺一文で十分な場合もあれば、前後段落や脚注本文を含めなければならない場合もある。そのためモデル上では開始・終了段落を記録し、抽出処理では参照表記とその周辺フィールドを一体として扱う。文脈の保存は検索精度の向上だけでなく、分類結果の説明可能性と学術的検証可能性を支える。

\subsection{碑文研究における参照意味論}

\subsubsection{カテゴリー設計}

碑文参照の意味タイプは、対象文献で観察される用法と碑文学上の基本的な研究行為をもとに定義した。設計時には、少数のサンプルをLLMによるドラフトで洗い出し、その結果を統合して確定した。これは確定的なオントロジーではなく、ケーススタディを通じて参照意味論を検証するための統制語彙である。

概念上は、いくつか次の群に整理できる。第一は史料の提示と識別であり、番号言及のみ（Citation only）と間版間同定（Cross-corpus identification）を含む。第二はテクスト提示であり、翻刻文提供（Transcription）と翻訳提供（Translation）を含む。第三は史料記述であり、内容説明（Content description）、年代提示（Dating）、出土地・所在記述（Provenance）、図版参照（Figure reference）を含む。第四は史料批判と分析であり、詳細分析（Detailed analysis）、テクスト校訂に関する分析（Textual criticism）、発掘調査・考古学的分析（Excavation / archaeological analysis）を含む。第五は研究史・編纂情報であり、発見者・刊行者言及（Discoverer / publisher）と目録化・分類学的分析（Cataloguing / taxonomic analysis）を含む。この階層化により、領域固有の細分類を維持しつつ、他の史料種へ一般化可能な上位意味論を与えられる。

なお、「番号言及のみ」は排他フラグであり、他のいずれのカテゴリーも付かない場合にのみコード側で自動付与される。したがって研究者が判断するのは実質的に十二のカテゴリーである。

\subsubsection{識別子と碑文実体}

碑文参照には、CIL、AE、ILTun、ILAfr、IRT、ILAlg、ILS、BAC、CRAIなどの略号と番号が用いられる。これらは碑文を指示する慣用的な識別子であるが、識別子そのものと碑文実体は区別しなければならない。同一碑文が複数の集成に収録される場合、CIL番号とAE番号は異なる文字列であっても、同一の碑文を指しうる。したがって、モデルではCorpus Referenceに集成名・巻・番号を保持しつつ、Inscription実体へ対応づける。

本ケーススタディでは、Inscriptionの外部識別子としてEpigraphik-Datenbank Clauss-Slaby（EDCS）のIDを利用する。これにより、文献中の参照表記が単なる文字列検索のキーとして残るのではなく、既存データベースの碑文レコードへ接続できることを狙う。EDCS-IDを碑文実体の主キーとすることで、異なる集成番号による複数の参照から、文献を横断して同一の碑文へ収束させられる。ただし、EDCSへの同定はすべての参照に対して成功するわけではなく、複数候補、未収録、表記揺れ、複合参照なども残る。ILSやBACのようにEDCSの識別子体系を持たない集成もあるため、この点からも、原表記・正規化表記・同定を別々に保存することが重要である。

\subsection{参照情報の抽出・分類・同定}

\subsubsection{対象資料}

実験では、OpenEdition上で公開されたローマ帝政アフリカ碑文学に関する章・論文を対象とした。具体的には、\textit{Antiquités africaines}、\textit{Pallas} などの学術誌論文と、Ausonius Éditions をはじめとする叢書所収の書籍章を含む計四十の文献を処理した。著者はAounallah、Christol、Guédon、Lefebvre、Leveau、Ricci らを含む。対象資料は機械可読なテクストと比較的明確な文書構造を持ち、碑文が議論の主要な根拠として用いられている。選定の目的は、特定の地域や研究テーマについて網羅的な引用索引を作ることではなく、一次史料参照の抽出、意味分類、外部碑文データベースへの同定まで、一貫して試行できる条件を確保することにある。

\subsubsection{三段階の処理パイプライン}

抽出処理は、参照抽出、意味分類、EDCS同定の三段階からなる。各段階は文献ごとの単一の正規ファイルに情報を追記し、処理途中で別の同一性をもつレコードが増殖しないようにする。書誌情報はCrossrefおよびOpenEditionのメタデータを用いて別途解決し、最終的にJSON、CSV、RDF/Turtleとして出力される。中間表現には学術標準のCSL-JSONを用いる。

第一段階では、碑文集成の略号と番号に対する正規表現を用いて参照表記を規則ベースで抽出する。規則ベース処理を先行させる理由は、比較的定型的な識別子の検出をLLMに全面的に委ねず、再現可能な候補集合を作るためである。抽出時には、参照文字列だけでなく、本文・脚注の別、段落位置、注番号、周辺テクストも保存する。

第二段階では、ローカル環境で実行する大規模言語モデルにより、各Mentionを固定語彙に基づいて多値分類する。プロンプトは、古代ローマ史および碑文学の専門家として文章を分析すること、適用可能なカテゴリーをすべて返すこと、本文に明示されない関係を推測しないこと、各カテゴリーの根拠となる短い引用を返すこと、JSON以外を出力しないこと、を要求する。表記揺れを防ぐため、出力はカテゴリー集合を再検査し、制約付きのコードで統制する。

第三段階では、正規化された参照キーをEDCSの碑文データベースと突合し、MentionをInscription実体へ接続する。この同定処理によって、複数の集成番号が同一史料に紐づけられ、一つの碑文に対する複数文献・複数箇所の言及を横断的に検索できる。

\subsubsection{LLM分類の位置づけ}

本手法においてLLMは、史料参照を自由に解釈して新たな知識を生成する主体ではなく、あらかじめ定義された語彙に基づいて文脈を分類し、その根拠を抽出する構成要素として位置づけられる。固定語彙は出力空間を制約し、JSON形式は後続処理との接続を容易にする。また、根拠箇所を同時に保存することで、専門家は分類結果だけでなく、その判断が依拠した文章を確認できる。

もっとも、カテゴリーの境界は常に明確ではない。内容説明と詳細分析の差異、単なる年代言及と年代比定の議論、翻刻された短い語句と本文引用の区別などは、専門家間でも判断が分かれうる。したがって、評価は単純な正解率に還元されず、分類語彙の定義、根拠箇所との整合性、研究上の検索要求に対する有用性、という複数の観点から行うべきである。

\subsection{抽出結果}

\subsubsection{参照数・文書位置・同定率}

四十の文献から、合計 2,796 件の碑文参照（Mention）を抽出した。EDCS-IDによる統合の結果、異なる碑文は 1,685 件であった。文書位置の内訳は、本文中に現れるものが 926 件、脚注中に現れるものが 1,870 件であった。参照の約三分の二が脚注に集中するこの偏りは、碑文参照の知識が本文だけでなく注の構造に強く依存することを示す。本文テクストのみを対象とした索引化では、参照の大部分を取り逃すことを意味する。

EDCS同定の内訳を表~\ref{tab:ch5-edcs}に示す。同定に成功したのは 2,054 件（突合成功 2,018 件、候補複数 36 件）であり、全体の約七割強にあたる。残りのうち 575 件は集成番号を持つが、該当レコードが見つからなかったもの、167 件はILSやBACのようにEDCSの識別子体系そのものを持たない集成への参照であった。完全自動化を前提とせず、未同定候補を専門家確認へ回す半自動的なワークフローが現実的である。

\begin{table}[htbp]
  \centering
  \caption{EDCS同定の内訳（全2,796件）}
  \label{tab:ch5-edcs}
  \begin{tabular}{lr}
    \hline
    状態 & 件数 \\
    \hline
    突合成功（matched） & 2,018 \\
    候補複数（ambiguous） & 36 \\
    該当なし（not found） & 575 \\
    識別子体系外（unmatchable, ILS/BAC等） & 167 \\
    \hline
    同定成功計 & 2,054 \\
    \hline
  \end{tabular}
\end{table}

\subsubsection{参照タイプの分布}

参照タイプの分布を表~\ref{tab:ch5-types}に示す。内容説明が 1,539 件と最も多く、間版間同定が 1,530 件、出土地・所在記述が 1,361 件で続く。詳細分析（1,259 件）、年代提示（1,121 件）、翻刻文提供（803 件）がさらに続く。碑文学の研究文献において内容記述・間版間同定・所在情報・分析・年代付与・翻刻が基礎的な役割を果たすことを量的に確認できる。一方、翻訳提供（149 件）、発掘調査・考古学的分析（133 件）は相対的に少ない。

\begin{table}[htbp]
  \centering
  \caption{参照タイプ別の出現件数（多ラベルを延べ）}
  \label{tab:ch5-types}
  \begin{tabular}{lr}
    \hline
    参照タイプ & 件数 \\
    \hline
    内容説明（Content description） & 1,539 \\
    間版間同定（Cross-corpus identification） & 1,530 \\
    出土地・所在記述（Provenance） & 1,361 \\
    詳細分析（Detailed analysis） & 1,259 \\
    年代提示（Dating） & 1,121 \\
    翻刻文提供（Transcription） & 803 \\
    発見者・刊行者言及（Discoverer / publisher） & 298 \\
    目録化・分類学的分析（Cataloguing / taxonomic analysis） & 262 \\
    テクスト校訂に関する分析（Textual criticism） & 256 \\
    番号言及のみ（Citation only） & 231 \\
    図版参照（Figure reference） & 204 \\
    翻訳提供（Translation） & 149 \\
    発掘調査・考古学的分析（Excavation / archaeological analysis） & 133 \\
    \hline
  \end{tabular}
\end{table}

件数の合計が総Mention数を超えるのは、一つのMentionが複数カテゴリーを持つ多値分類だからである。この点は、参照を一つの目的へ還元しないというモデル上の判断を裏付ける。たとえば翻刻と間版間同定が同一脚注に併記され、そこへ翻訳や年代情報が加わることがある。意味的接続は、排他的な文献分類ではなく、狙いの言及が持つ複数の学術的機能を記述する必要がある。

\subsection{意味的参照ネットワークの構築}

\subsubsection{RDF知識グラフへの変換}

抽出結果はRDFへ変換し、文献・文書構造・Mention・Corpus Reference・Inscription・Relation Assertionをノードとして接続する。中心となるMentionから、\texttt{refersTo}によってInscriptionへ、\texttt{hasRelation}によってRelation Assertionへ、\texttt{hasCorpusReference}によって集成参照へ、文書位置プロパティによってParagraphまたはFootnoteへ接続する。文書側ではParagraphとFootnoteがSource DocumentまたはBook Sectionへ階層化される。MentionはWeb Annotation語彙に準拠し、InscriptionはEDCSレコードへ \texttt{owl:sameAs} でリンクされる。生成したグラフは、四十の文献を同一のTurtleに束ね、碑文をEDCS-IDで横断統合したもので、2,796件のMentionノードと1,685件のInscriptionノードを含む。

この構造の重要な点は、同一碑文への複数のMentionを一つの史料実体の周囲に集約できることである。従来の参照索引が「史料番号→文献」の対応表であったのに対し、知識グラフは「史料→言及→文脈→参照機能→文献」という経路を提供する。利用者は、ある碑文を翻刻する文献、年代を論じる文献、異読を提示する文献、出土状況を扱う文献を、意味タイプで絞り込んで取得できる。

\subsubsection{共参照と文献横断接続}

同じ碑文を参照するMention同士は、共通のInscriptionノードを介して文献横断的に接続される。さらに、同一箇所で共に引用される碑文をco-cited-withとして記述すれば、研究文献上で同じ議論に召喚される史料群を探索できる。これは、単純なキーワード共起ではなく、正規化・同定された史料実体の共参照ネットワークである。

このネットワークは、研究書の章立てやOpenEdition上のページ構造を超えて、碑文を軸とした再編成を可能にする。個別章に分散した議論を、一つの研究対象のもとに集約でき、集成ノードと碑文実体の双方へ文書セクションを位置づけることで、横断的な閲覧を実現する。

\subsection{検索および生成AIへの応用}

\subsubsection{意味タイプを利用した史料中心検索}

知識グラフ化によって、一次史料から二次文献への探索は、単なる全文検索よりも精密になる。利用者は、「その碑文を引用する文献」だけでなく、「その碑文の年代を提示する箇所」「本文を翻刻する脚注」「異なるreadingを論じる箇所」などを指定できる。検索結果には、文献の書誌情報、該当言及位置、意味タイプ、根拠テクスト、原文への深いリンクが付与でき、結果から原文脈へ戻る経路が確保される。

\subsubsection{二つのコンテクストを統合するGraphRAG}

意味的接続モデルの応用例として、この知識グラフをGraphRAG型の検索システムに組み込むことができる。第一の情報流はベクトル検索であり、質問を埋め込み表現へ変換し、該当する碑文テクストとメタデータを取得する。第二の情報流はグラフ検索であり、グラフ検索クエリによって当該碑文を参照する研究文献、著者やタイプ、意味タイプ、深いリンクを取得する。二つの情報流を、碑文コンテクストと近現代の研究文献コンテクストとして統合し、言語モデルへ渡す。

これにより、生成される応答は、一次史料の内容だけでなく、その史料に対する研究上の翻刻・年代・解釈・批判を参照できる。ベクトル検索は非構造化テクストの柔軟な検索を担い、知識グラフは同定済みの参照関係と意味タイプを担うという役割分担である。ただしGraphRAGは本章の目的そのものではなく、意味的接続モデルを可能にする応用例として位置づけるべきである。本章の主要な貢献は、生成応答の性能ではなく、一次史料と研究文献のあいだに検証可能な意味的経路を構築する点にあり、生成AIはその経路を利用する一つのインターフェースである。

\subsection{課題と限界}

\subsubsection{デジタル化とアクセス可能性}

第一の課題は、二次文献のデジタル化状況である。多くの研究文献はPDFまたは紙媒体でのみ利用できる。PDFも論理構造を持つ機械可読データとは限らない。OCR誤り、脚注と本文の分離、段組、図版、特殊文字は抽出精度に影響する。また著作権のため全期を集約・再配布できない場合が多い。したがって、意味的索引では、必要最小限の文脈と根拠表現を保存しつつ、可能な限り原公開元へのリンクを維持する設計が必要である。

\subsubsection{参照形式と同定の多様性}

第二の課題は、参照正書法の多様性である。略号、巻号、番号範囲、旧型・新型の対応、ibid.型の省略、複数史料の一括参照は、集成・言語・時代・研究者によって異なる。今回の正規表現は対象資料に現れる主要な集成を扱うが、別のコーパスへ適用する際には規則の追加と検証が必要となる。規則ベース抽出と外部データベース突合を分離するのは、この拡張性を確保するためでもある。

\subsubsection{意味分類の評価}

第三の課題は、意味分類の評価である。参照カテゴリーは自然種ではなく、研究目的に応じて設定された分析概念である。カテゴリー間には重なりがあり、文脈範囲の選択によって判定が変わりうる。そのため、単一のgold standardとの一致率だけでは十分ではなく、複数専門家による一致度、誤分類の質的分析、根拠箇所の適切性、検索タスクへの寄与、を組み合わせて評価する必要がある。

また、LLMのモデル更新やプロンプト変更によって結果が変動する可能性もある。分類結果をRelation Assertionとして独立させ、使用モデル・プロンプト・実行時刻・根拠箇所・検証状態を来歴情報として保持すれば、再現性と更新可能性を高められる。これはAI出力を確実的知識として直接登録するのではなく、検証対象となる解釈提案として扱う方法である。

\subsubsection{カテゴリー体系の拡張}

今回のカテゴリー体系は碑文学の基本的な参照機能を対象とするが、すべてを網羅してはいない。たとえば、特定の歴史命題を史料がsupportまたはcontradictする関係、先行刊本のreadingを採用・拒否する関係、碑文の真正性を論じる関係などは、より細かなモデル化が可能である。他方、過度に細分化すれば分類の一貫性と他分野への移植性が損なわれる。汎用的上位概念と碑文学固有の下位概念を分け、実際の検索要求に基づいて段階的に拡張することが望ましい。なお、統制語彙の整備は現在も進行中であり、分類実装で用いる語彙とオントロジー上の統制語彙を一致させる作業を伴う。

\subsection{他の一次史料への一般化}

本章の実験は碑文を対象としたが、Mentionを中心とする構造自体は他の一次史料にも適用できる。文献史料では刊本・巻・章節・行番号、パピルスではコレクション番号と再編集、貨幣ではカタログ番号と刻印、美術作品では作品識別子と図像解釈、考古資料では出土コンテクストと取扱分類が、それぞれCorpus ReferenceとRelation Typeに相当する。

一般化のためには、二層の語彙設計が有効である。上位層には、提示・識別、内容提示、記述、分析・批判、研究史・編纂といった領域横断的な参照機能を置く。下位層には、翻刻、間版間同定、欠損補塡、取扱分類など、史料種ごとの研究実践を置く。これにより、異なる資料領域の参照を比較可能にしつつ、各領域に必要な精度を保持できる。

さらに、一次史料参照が研究行為として捉えられるならば、二次文献は単に史料への入口ではなく、史料に対する解釈が蓄積される場である。この蓄積を構造化することは、史料データベースと書誌データベースを接続するだけではなく、史料から学術的主張が形成される過程を機械可読化することを意味する。したがって、意味的接続モデルは、歴史知識グラフに研究史と解釈の層を組み込むための基礎となる。

\section{一次史料に対するセマンティック・アノテーション}

\subsection{問題設定：史料側からの意味的接続と対象媒体の拡張}

5.3で扱ってきたのは、二次文献に既に埋め込まれた一次史料への参照を抽出し、Mentionを軸とする構造化を通じて意味的接続を可能にするアプローチであった。この方向性は、既存の研究蓄積が豊富な領域においては大きな成果を上げる一方で、二次文献における言及の網羅性に強く依存する。すなわち、二次文献において明示的に言及されていない史料や、そもそも既存の研究蓄積が薄い新出資料・未刊行資料については、この抽出的アプローチだけでは意味的接続のネットワークに組み込むことができない。

そこで本節では、逆方向のアプローチとして、一次史料の側から関連する二次情報を能動的に紐づけていく枠組みを検討する。すなわち、史料の特定箇所に対して、それを解釈する上で必要となる書誌情報・典拠情報・関連メディアを、研究者自身が付与していくというアノテーション実践の設計である。この二つの方向性は排他的な関係にあるのではなく、最終的には同一の知識グラフ上で合流し、相互に補完する関係にある。この統合像については本節末尾で改めて論じる。

ただし、5.3から5.4への移行は、単に方向性の反転にとどまらない。ここには対象媒体の質的な拡張が伴う。5.3で展開した抽出的アプローチにおいて対象となるのは、二次文献のテキスト内に現れる史料への言及そのものであり、参照する側の二次文献がテキストであることが決定的な条件であった。この条件のもとでは、参照される一次史料の物理的形態は、抽出作業にとって本質的な問題とならない。たとえば碑文の場合、二次文献はその翻刻文を掲げたり、年代を論じたり、出土地を記述したりするが、これらはいずれも二次文献のテキスト平面内で完結する参照行為であり、碑文そのものが石板・青銅板・パピルスのいずれであるかは、参照抽出のプロセスにおいては直接には現れてこない。Mention中心の設計が有効に機能したのは、参照の起点がテキスト平面上に配置され、参照される対象は識別子と記述を介して間接的に指示されるだけで十分であったからである。

しかし、史料側からの意味的接続を考えるとき、状況は根本的に異なる。ここでは参照の起点が一次史料そのものに置かれるため、史料の物理的形態が正面から問題となる。文化遺産研究において解釈の対象となる一次史料は、碑文や写本のようにテキストを主要な構成要素とする媒体だけでなく、絵画・彫刻・工芸品・建築・出土遺物・古地図といった視覚情報を主体とする媒体、さらには近年整備が進む3Dスキャンデータや構造化復元モデルにまで及ぶ。これらの非テキスト媒体においては、史料そのものが可読なテキストを含まない以上、5.3のような「本文からの参照抽出」は原理的に成立しない。史料のどの部分に対する解釈なのかを研究者が明示的に指定し、その上に関連情報を紐づけていくほかない。この点で、非テキスト媒体を対象とする限り、意味的接続は必然的に史料側からの能動的なアノテーション行為として立ち現れる。

さらに、この視点はテキストを主要な構成要素とする史料に対しても新たな含意を持つ。碑文の石材そのものの摩耗状況、装飾要素の様式、彫刻面の物理的な配置といった、テキストに還元されない側面についての解釈は、二次文献のテキスト平面内で言及されうるとしても、その史料の視覚的・物理的な表現を対象として部分領域を明示的に指定した上で付与されるべき性格のものである。したがって、たとえテキスト史料であっても、その物理的・視覚的側面を扱う限りにおいては、非テキスト媒体と同じく史料側からの能動的アノテーションとして扱われることになる。逆に言えば、画像・3Dを対象として設計される枠組みは、テキスト史料の物理的側面をも包含しうるより一般的な意味的接続の基盤となりうる。実際、後述する本研究のシステムは、2D画像・3Dモデルとあわせてtextualエディタを備えており、TEIで構造化されたテキスト資料の特定箇所に対しても同一の枠組みでアノテーションを付与できる。この統合可能性については本節末尾で改めて論じる。

以上の理由から、本節ではまず対象媒体として、時間性を含まない非テキスト媒体、すなわち2D画像資料および3Dモデルを中心に据える。時間性を含む音声・動画・アニメーション付き3Dは、時間性を含まない資料の枠組みを拡張することで扱えるため、まず基盤となる枠組みの設計を優先する。データ表現の基礎としてはIIIF Presentation API 3.0を採用し、これによって扱われる2D画像・3Dモデルのマニフェストを一次史料の入力とする。5.3で扱ったローマ帝政アフリカの碑文についても、そのかなりの部分がすでに3Dスキャンデータとして整備されつつあり、テキスト史料としての碑文（5.3で対象とした側面）と、物理的・視覚的対象としての碑文（本節で対象とする側面）は、同一の史料の二つの相として同一の知識グラフ上に位置づけられる。この点でも、章内における対象の連続性は具体的に担保される。

以下、本節では、この能動的アノテーションを実現するための概念モデルを段階的に提示する。まず5.4.2では、複数主体による重層的な解釈の集約を可能にする「領域リソース」概念を導入し、5.4.3ではその実装的基礎としてIIIF Presentation APIの拡張を論じる。5.4.4以降では、領域リソース上に付与されるアノテーションの内実、すなわち書誌・典拠・メディアという三種のLinked Resourceの意味論的枠組み（CHAO v0.1）を展開する。

\subsection{共同的アノテーションと「領域リソース」の導入}

一次史料に対する能動的アノテーションを歴史研究の実践として捉え直すとき、まず問題となるのは、単一の研究者による孤立したアノテーションではなく、複数の主体が同一史料の同一箇所に対して重層的な解釈を積み重ねる過程をいかに記述するか、という点である。歴史研究における解釈は本質的に多声的であり、ある碑文の装飾要素について異なる様式帰属を主張する研究者、同じ人物名の同定について異なる候補を挙げる研究者、同じ場面の意味について異なる読み解きを提示する研究者が、同一史料の同一箇所を対象として並立する。この多声性を意味的接続の枠組みに取り込むことが、5.4節の設計上の中心的な要件となる。

ところが、既存のアノテーションモデルの多くは、この要件を構造的に満たしていない。従来のアノテーション事例においては、アノテーションと注釈領域は一対一で対応することが前提とされ、注釈領域に関する座標情報は個々のアノテーションに内包する形で記述されてきた。この設計のもとでは、仮に複数の研究者が資料上の「同一の」箇所について解釈を付与する場合であっても、個々のアノテーションは互いに独立したデータとして生成される。座標情報が各アノテーションに埋め込まれる以上、たとえ人間の目には同一領域を指すように見えても、データ構造上は完全に別個の対象に対して別々のアノテーションが付与されているという扱いになる。マルチメディア資料においては、複数のユーザーがマウス操作等で座標を選択する際に完全な同一性を確保することは実質的に困難であり、この課題はいっそう深刻となる。

このような構造は、意味的接続の観点からみて根本的な欠落を持つ。同一領域に対する複数の解釈が構造的に集約されえない以上、その領域について蓄積された解釈の全体像を検索・比較することも、解釈相互の関係（支持・反論・補完）を記述することも不可能となる。多声的解釈の記録と共有は、単に複数のアノテーションを並列させることではなく、同一の対象を指しているという事実を機械可読な形で保証することを要求するのである。

本節が採用する解決策は、資料上の物理的領域を、個々のアノテーションから独立した「領域リソース」として定義することである。領域リソースとは、アノテーション対象オブジェクト上の特定の物理的領域を抽象化したデータリソースであり、それ自体では解釈に関わる意味内容を含まない。図~\ref{fig:ch5-region-resource}に示すように、領域リソースの導入は、これまで不可分に記述されてきた注釈領域情報と解釈情報を明確に分離する機能を有する。異なる主体によって生成された複数のアノテーションが、共通の一つの領域リソースに対して付与されることで、注釈領域については複数の異なるユーザーで共有しつつ、その上にそれぞれの解釈を重ねていくことが可能になる。

\begin{figure}[htbp]
  \centering
  % \includegraphics[width=0.85\linewidth]{figures/chapter5/region-resource.png}
  \caption{「領域リソース」を導入したアノテーションの概念図}
  \label{fig:ch5-region-resource}
\end{figure}

この設計は、5.3で採用した設計原理と本質的に同型である。5.3では、参照表記を担うCorpus Referenceと、それが指示する史料実体としてのInscriptionを分離した。これによって、CIL番号とAE番号という異なる表記が同一の碑文実体を指す事実を構造的に表現できるようになった。5.4における領域リソースの導入は、この分離原理を史料側のアノテーション実践に対して適用したものと理解できる。すなわち、複数の異なる主体がそれぞれ生成する座標情報付きアノテーション（Corpus Referenceに相当）と、それらが指示する共通の物理的領域（Inscriptionに相当）を区別することによって、多主体的解釈の構造的集約を可能にする。

領域リソースは個々のアノテーションにおける解釈からは独立したデータ資源であるため、複数のユーザー、あるいはユーザーが属するプロジェクトを跨いで共有することができる。この点で、領域リソースはいわば種々の個人やプロジェクトの共有財産として機能する。特定のプロジェクトが生成した領域リソースに対して、別のプロジェクトが独自のアノテーションを付与することが可能となり、オブジェクトの特定領域についての多様な解釈を集約・蓄積することが可能になる。ここでは、「共有される事実としての物理的領域」と「解釈としてのアノテーション」という二層構造が構造的に区別されており、前者はプロジェクト横断で共有され、後者は個々の主体・プロジェクトの解釈行為として帰属先を保持する。この二層構造は、歴史研究における解釈の主体性と共有性の緊張関係を、データ構造の水準で反映するものである。

\subsection{IIIF Presentation APIの拡張}

領域リソースを軸とする共同的アノテーションモデルを実際にデータとして記述・共有するためには、そもそもアノテーションの対象となる資料そのものがウェブ上で共有されている必要がある。すなわち、特定のデータ提供システムに依存するのではなく、共通の仕方でどこからでもアクセス可能な環境を整える必要がある。こうした要件を実現するためにデジタル人文学の分野で一般的に用いられるのが、IIIF（International Image Interoperability Framework）である。国際的な画像の相互運用のための枠組みであるIIIFは、Image API、Presentation APIなど複数のAPIを通して画像データおよびメタデータの簡易な公開・共有・操作を実現し、国内外の多くの文化機関において採用されている。二次元の画像・動画のアノテーションについてはPresentation API 3.0ですでに定義されており、研究実践も存在する。他方、3Dについては未対応であったが、近年、3D対応を議論するワーキンググループが規格の策定を進めており、Presentation API 4.0から本格導入される予定である。試験的な形での3D対応はすでに現在でも行われており、Universal Viewer等では3Dデータを含むマニフェストの読み込み・表示がすでに可能になっている。

こうした状況を踏まえ、本研究でもアノテーション記述にあたっては基本的にIIIFに準拠する方針を採る。IIIFのPresentation API 3.0では、アノテーションは \texttt{type} を \texttt{Annotation} とするJSONオブジェクトとして記述され、\texttt{target} によってCanvas全体またはその一部分を参照する。Canvasの部分領域を対象とする場合、\texttt{target} には \texttt{type} を \texttt{SpecificResource}（SR）とするオブジェクトを用い、\texttt{source} に対象となるCanvasのURIを、\texttt{selector} に領域を特定するための座標情報を記述する。2D画像の矩形領域はFragmentSelectorのvalueに \texttt{xywh=248,382,62,73} の形式で指定できる。一方、3Dモデル上の点や面の指定には、3Dデータ向けに定義されたPointSelectorやPolygonZSelectorなどのSelectorを用いる。

ただし、5.4.2で導入した領域リソースをIIIF上で実装するにあたっては、既存のIIIFの記述を一部拡張する必要がある。既存のIIIF Presentation APIにおいては、資料の部分領域はSRによって記述されるが、SRは各アノテーションのtargetにおいて対象領域を特定するための補助的な記述要素として位置付けられており、その意味論はWeb Annotation Data Modelなどの外部仕様に委ねられている。Presentation API自体も表示（presentation）を目的とした仕様であるため、SRを複数のアノテーションから共有される独立した知識リソースとして永続的に識別・再利用するための仕組みは規定されていない。SRに \texttt{id} を付与すること自体は可能であるものの、その \texttt{id} は当該アノテーション内の対象を識別するためのものであり、同一領域を表す永続的なリソースとして共有・再利用されることは想定されていない。そのため、実運用上は、SRは個々のアノテーションに従属する一時的な領域指定として扱われ、同一対象に対する複数のアノテーション間で共通の参照対象となることはない。

これに対して本研究では、図~\ref{fig:ch5-iiif-shared-sr}に示すように、SRを領域リソースとして独立に定義し、永続的なURIを付与する。同一領域を対象とする複数のアノテーションは、この共通のURIを持つSRを \texttt{target} として記述することで、既存のIIIFビューアとの互換性を維持しながら、同一領域を共有するアノテーション群として識別・集約できるようにする。各アノテーションには、既存のIIIF仕様に従って \texttt{source} および \texttt{selector} が保持されるため、既存のビューアにおいても従来通り表示可能である。一方で、SRに付与された永続的なURIを共通の識別子として利用することで、同一対象に対する複数の解釈や注釈を領域単位で管理できるようになる。さらに、このURIは異なるマニフェストやシステム間においても共通の参照対象として利用可能であり、IIIFとの相互運用性を維持しながら、同一対象に対する複数解釈の共存と領域単位での知識管理を実現する。

\begin{figure}[htbp]
  \centering
  % \includegraphics[width=0.95\linewidth]{figures/chapter5/iiif-shared-sr.png}
  \caption{同一「領域リソース」に属する二つのアノテーション記述（IIIF出力例）}
  \label{fig:ch5-iiif-shared-sr}
\end{figure}

この拡張は、5.3において採用した「参照表記の冗長保持による下位互換性の担保」という設計判断とも通底している。5.3では、正規化された参照キーとEDCS-IDによる同定を導入しつつ、原表記（raw reference）を捨てずに保持することで、既存の書誌検索との整合性を保った。5.4におけるSRの拡張も同様に、\texttt{source}と\texttt{selector}を各アノテーション内に完全展開したまま、\texttt{id}による集約の層を追加する。既存のビューアはこの \texttt{id} を単にアノテーション内の識別子として無視すればよく、追加された集約情報を理解するシステムだけが、共通URIを介した領域単位の管理を実行する。この二層構造によって、意味的接続の精緻化と既存インフラとの後方互換性が同時に達成される。

\subsection{アノテーションの概念モデル：CHAO v0.1におけるLinked Resource}

領域リソースが「どこに」対する解釈かを担うのに対し、その上に付与されるアノテーションは「何を」参照し、「どのような」解釈を提示するかを担う。本研究では、この参照される二次情報の側を三種類のリソースに分類する枠組みとして、Cultural Heritage Annotation Ontology（CHAO v0.1）を採用する。CHAO v0.1は、文化遺産資料に対するアノテーションにおいて参照されるリソースを、書誌（BibliographicResource）、典拠（AuthorityResource）、メディア（MediaResource）の三種に分類し、それぞれに下位区分を与える。

書誌（BibliographicResource）は、当該史料を論じたり参照したりする学術的著作を指す。CHAO v0.1では書誌をPrimary（一次資料的書誌、たとえば同時代の史料集）、Secondary（二次研究文献）、Report（発掘報告や調査報告）の三種に細分化する。この分類は、5.3で扱った参照抽出の対象と直接に対応する。5.3において二次文献から抽出されたMentionが指示する側の情報がInscription実体であったとすれば、5.4において史料側から付与される書誌参照が指示する側の情報がBibliographicResourceである。両者は方向を逆にした同一の関係であり、共通の書誌URI（Crossref DOI、OpenEditionのハンドル等）を介して知識グラフ上で結合される。

典拠（AuthorityResource）は、史料の解釈において参照される権威データ上のエンティティを指す。CHAO v0.1ではこれをPerson、Place、Event、Object、Period、Region、Culture、Language、Script、Conceptといった下位区分に整理する。この分類は、4章で導入されたHIMIKOモデルにおけるエンティティ分類と対応関係を持ち、Wikidataや各種の権威ファイル（VIAF、Getty AAT、Pleiades等）が付与するQIDや識別子を通じて外部と接続される。史料上の特定領域に対して、そこに描かれた人物、そこに言及される場所、それが行われた時代といった典拠エンティティを紐づけることで、当該領域が指し示す指示対象の同定が明示される。

メディア（MediaResource）は、当該史料を補完する視覚的・文書的資料を指す。CHAO v0.1ではImage（File・IIIF）、Video（YouTube等）、3D Modelの下位区分を設ける。史料上の特定領域について、その領域を別角度から撮影した写真、関連する場面の映像、比較対象となる別史料の3Dモデル等を紐づけることが可能となる。特にIIIFで表現されたメディアは、それ自体が別の一次史料としてこのシステムに登録されうるため、メディア参照を介した史料間の相互参照ネットワークが構築される。

この三分類の意義は、5.3のMention中心の抽出モデルとの対比において明確になる。5.3では、二次文献のテキスト内に匿名的に埋め込まれた参照を、Mentionという中間ノードを介して構造化した。ここでは参照される側（Inscription）は既に確定した史料実体である一方、参照する側（Publication）が一つの文献という粒度で提示されるため、参照の対象を細分化することが主要な課題となった。これに対して5.4では、参照する側（史料の特定領域）は領域リソースによって既に明示的に位置づけられる一方、参照される側（書誌・典拠・メディア）が多様な種類に及ぶため、これを型付けて構造化することが主要な課題となる。二方向のアプローチは、Mention単位と領域リソース単位という異なる中間ノードを持ちながら、いずれも「参照する側と参照される側の中間に構造化ノードを挿入する」という共通の設計原理に立つ。

\subsection{関係性の意味論：Relation Hierarchy}

Linked Resourceの三分類は、参照される側の類型論を与える一方で、参照が果たす機能そのものは記述しない。史料上の特定領域とそこに紐づけられる書誌・典拠・メディアの間には、単に「関連する」というだけでは捉えられない多様な意味的関係が存在する。ある書誌はその領域を単に言及するのみである一方、別の書誌はその領域を詳細に分析する。ある典拠はその領域が描く人物そのものであり、別の典拠はその領域を理解するために比較対照される事例である。この関係の型付けを担うのが、本研究が採用するRelation Hierarchyである。

Relation Hierarchyは、\texttt{public/ontology/relation-hierarchy.json} に定義される二階層の統制語彙として実装される。上位分類はDirect RelationとConceptual Relationの二種からなる。Direct Relationは、史料領域と参照対象の間の直接的・字義的な関係を記述し、\texttt{mentions}、\texttt{depicts}、\texttt{illustrates}系のプロパティを含む。これらは書誌・典拠・メディアの各リソース種別に対して個別に語彙が用意されている。Conceptual Relationは、より解釈的・文脈的な関係を記述し、\texttt{contextualizes}（およびその四つの下位分類）、\texttt{compares\_with}、\texttt{provides\_typology}、\texttt{associated\_with}、\texttt{classified\_as}などを含む。

この二階層の設計は、5.3のRelation Assertionにおける参照タイプ体系と密接な対応関係を持つ。5.3の参照タイプ（内容説明、間版間同定、翻刻文提供、詳細分析等）は、いずれも「二次文献が一次史料に対してどのような機能を果たすか」を記述する語彙であった。5.4のRelation Hierarchyは、これを二次情報全般（書誌・典拠・メディア）に対する史料側からの関係記述に拡張したものと理解できる。実際、5.3の「内容説明」は\texttt{describes}系、「詳細分析」は\texttt{analyzes}系、「翻刻文提供」は\texttt{transcribes}系というように、Relation HierarchyのDirect Relation下位に対応する。この対応関係は、二方向のアプローチによって蓄積される解釈が、同一の意味論的空間の中で比較・集約可能となることを保証する（この点は5.4.9で改めて論じる）。

Relation Hierarchyの重要な特性は、リソース三分類（書誌・典拠・メディア）と関係語彙が直交する点にある。同じ\texttt{mentions}関係が書誌に対しても典拠に対しても適用可能であり、同じ\texttt{contextualizes}関係が書誌・典拠・メディアのいずれに対しても適用可能である。この直交性によって、リソースの種類に応じて別々の語彙を覚える必要がなく、意味論的な統一性が保たれる。同時に、リソース種別ごとに固有の関係（たとえば書誌に対する\texttt{cites}、典拠に対する\texttt{depicts}）も併存させることで、領域固有の精度も担保される。この二層設計は、5.3で議論した「汎用的上位概念と領域固有の下位概念の共存」という一般化戦略と同型である。

なお、5.3のCiTOに関する議論で触れたように、関係語彙の重要性は、単なる有向辺としての参照リンクを、意味づけを伴う参照関係へと拡張する点にある。Relation Hierarchyの導入によって、史料側から付与されるアノテーションは、単に「この領域はこの書誌に関連する」という平坦な記述ではなく、「この領域についてこの書誌がこのような機能を果たしている」という機能的記述として展開される。この記述性は、後述する検索・比較・再利用の精度を根本的に規定する。

\subsection{付与行為の来歴：E13 Attribute Assignment}

領域リソース、Linked Resource分類、Relation Hierarchyという三つの層は、アノテーションの構造的枠組みを与えるが、これだけではアノテーションが誰によっていつ何を根拠に付与されたかという来歴情報を扱わない。歴史研究における解釈は本質的に主体的な行為であり、その解釈を提示した研究者、提示された時点、提示に伴う根拠づけといった付与行為の情報は、解釈の内容そのものと同等に重要である。この来歴情報を構造的に扱う枠組みとして、本研究ではCIDOC CRMのE13 Attribute Assignmentクラスを採用する。

E13 Attribute Assignmentは、あるエンティティに対して何らかの属性・関係が付与される行為そのものを一つのエンティティとして記述する。これによって、付与された属性・関係の内容だけでなく、その付与を実行した主体（\texttt{addedBy}）、実行時点（\texttt{addedAt}）、付与の際のコメントや根拠（\texttt{addedComment}）といったメタ情報を、独立したノードとして保持できる。本研究のシステムにおいては、書誌・典拠・メディアの各リソース紐づけが、それぞれ独立したE13 Attribute Assignmentとして記録される。あるアノテーションが特定の書誌を参照する場合、その参照は「アノテーションが書誌を参照する」という二項関係ではなく、「研究者Aがある時点にコメントBとともに、アノテーションと書誌の間の\texttt{cites}関係を宣言した」というE13行為として記述される。

この設計は、5.3の末尾で議論した「LLM分類結果をRelation Assertionとして独立ノード化する」という方針と本質的に同型である。5.3では、LLMによって付与された意味分類を確定的属性値として保存するのではなく、使用モデル・プロンプト・実行時刻・根拠箇所・検証状態を来歴情報として保持することで、再現性と更新可能性を確保した。5.4のE13 Attribute Assignmentは、この方針を人間による解釈行為へと拡張したものと理解できる。両者は共に、歴史知識を確定的事実の集積としてではなく、検証可能な解釈行為の履歴として記述するという本論文全体の方針を実装している。

E13的な来歴管理は、複数主体による解釈の相互関係を記述する上でも本質的な機能を担う。5.4.2で述べたように、領域リソースの共有可能性は同一領域に対する複数の解釈の並立を可能にするが、これらの解釈は単に並立するだけでなく、しばしば相互に支持・反論・補完の関係を持つ。本研究のシステムでは、同一領域ノードに紐づく複数のアノテーション間に、\texttt{supports}、\texttt{challenges}、\texttt{supplements}という三種の関係を付与できる。これらは Web Annotation Data Modelの \texttt{oa:motivatedBy} の下位語彙として位置づけられ、RDF出力においては \texttt{rdf:Statement} によるリフィケーションを通じて、関係そのものに対する付与者・日時・コメントを記述できる。すなわち、「アノテーションAはアノテーションBを支持する」という関係もまた、それを主張した主体・時点・根拠を伴った一つの解釈行為として記録される。

この重層的な来歴管理によって、同一領域に対する解釈の集積は、単なる解釈の並列的な列挙ではなく、解釈者間の対話の記録として構造化される。ある研究者Aの解釈に対して、研究者Bがそれを補完する新たな解釈を提示し、研究者Cが両者のいずれをも部分的に否定する別解釈を提示するといった議論の展開を、意味的接続の枠組みの中で機械可読な形で保持することが可能となる。歴史研究の解釈が本質的に対話的・累積的な営為であることを踏まえるとき、この対話の構造化は、5.3では扱いえなかった能動的アノテーションの側面の中でもとりわけ重要な機能を担う。

\subsection{記述の場としての本文：Markdown記法とリソース内参照}

これまで論じてきた領域リソース、Linked Resource分類、Relation Hierarchy、E13 Attribute Assignmentは、アノテーションを構造化するための枠組みであり、いずれもアノテーションの構造的側面を扱う。しかし、実際のアノテーションにおいて研究者が最も労力を投じるのは、これらの構造的要素を組み合わせながら、自然言語による記述としての解釈本文（Description）を執筆する作業である。本節では、この本文執筆を支える記法とその意味論を論じる。

本研究のシステムでは、アノテーション本文をCommonMark + GFM（GitHub Flavored Markdown）準拠のMarkdown文字列として記述する。Markdownの採用は、記法の簡便さと表現力の両立を意図したものである。研究者は特別なエディタや専門的な記法を習得することなく、見出し・箇条書き・強調・引用といった基本的な文書構造を用いて解釈を記述できる。一方で、Markdownはプレーンテキストとして保存されるため、機械可読性を損なわない。

本研究の記法上の独自性は、Markdownの標準的なリンク記法を拡張して、同一アノテーションに紐づいたLinked Resourceへの内部参照を可能にした点にある。すなわち、本文中で \texttt{[label](\#bib/<id>)} と記述すれば、当該アノテーションに紐づく書誌エンティティ（BibliographicResource）のうち指定IDのものが参照される。同様に \texttt{[label](\#auth/<id>)} は典拠エンティティ、\texttt{![caption](\#media/<id>)} はメディアエンティティを参照する。この記法によって、研究者は本文中で「〜については[A氏の研究](\#bib/xyz)によれば」「[この地域](\#auth/pqr)における」「[比較図版](\#media/uvw)を参照」といった形で、Linked Resourceを本文の論述の流れの中に自然に埋め込むことができる。

この記法の意義は、本文と紐づけリソースの分離を回避し、両者を統合的に扱える点にある。従来のアノテーションシステムでは、本文はプレーンテキストとして記述され、参照リソースはそれとは別のフィールドで管理されるという分離が一般的であった。この分離は、本文中のどの記述がどのリソースに対応するかという対応関係を機械可読な形で保持することを困難にする。本研究のMarkdown内参照記法によって、本文中の各記述と紐づけリソースの対応が明示的に記録される。

さらに、この内部参照記法は、5.3で参照したCiTO（Citation Typing Ontology）の実装例としても機能する。本研究のシステムでは、アノテーション本文のMarkdownを走査し、\texttt{\#bib/<id>} 形式の参照を \texttt{cito:cites} 関係として、他の意味論的な参照パターンを \texttt{cito:discusses} 関係として自動的に抽出し、RDF出力に含める。すなわち、研究者はMarkdown記法の自然な延長として書誌参照を執筆するだけで、CiTOに準拠した意味づけ引用ネットワークが自動的に生成される。5.3が二次文献の既存テキストからLLMによって参照タイプを抽出したのに対し、5.4では研究者が本文執筆の過程で意味的接続を明示的に宣言することになる。両者は共に、「参照関係に意味づけを与える」というCiTOの発想を、それぞれ抽出的・宣言的な形で実装している。

\subsection{システム実装と出力}

以上の概念的枠組みは、\texttt{IIIF Semantic Editor} と名付けられたWebベースのアノテーションシステムとして実装されている。本節では、システムの構成、ユーザーインターフェース設計、およびIIIF/RDF二形式による出力について概観する。

システムはNext.js（App Router）をフレームワークとし、Firestoreをデータストアとして用いる。データは、アノテーション本体を保存する \texttt{test} コレクション、領域ノードの座標情報を保存する \texttt{regions} コレクション、オブジェクトメタデータを保存する \texttt{manifest\_metadata} コレクションに分散して管理される。\texttt{test} コレクションと \texttt{regions} コレクションの分離は、5.4.2で論じた「解釈情報と領域情報の分離」という設計原理を、データストアの物理レベルで実装したものである。なお、アノテーションの帰属と編集権限は研究プロジェクト単位で管理される仕組みを備えているが、その詳細は本論の主題を外れるため割愛する。

システムは2Dエディタ、3Dエディタ、textualエディタの三つのモードを備え、それぞれが対応する媒体（2D画像、3Dモデル、TEIテキスト）のアノテーションを支援する。ユーザーはIIIFマニフェスト入力欄にURLを入力することで対象史料を読み込み、インタラクティブな操作を通じてアノテーションを作成・編集・共有できる。図~\ref{fig:ch5-editor}に、3Dエディタの編集画面を示す。

\begin{figure}[htbp]
  \centering
  % \includegraphics[width=0.95\linewidth]{figures/chapter5/editor-screen.png}
  \caption{IIIF Semantic Editorの編集画面}
  \label{fig:ch5-editor}
\end{figure}

インターフェース上では、ユーザーがアノテーション対象の3Dモデルや画像上でダブルクリックまたはポリゴン選択を行うと、まず領域リソースが作成され、続いてその領域に紐づく最初のアノテーションの入力が促される。一方、既存の領域リソースのマーカーをクリックした場合は、その領域に紐づくアノテーションの一覧が右パネルに表示され、既存のアノテーションを選択して参照・編集するか、新規アノテーションを追加するかを選択できる。多くの既存アノテーションシステムにおいては、領域選択から直接アノテーション入力に進むが、本システムでは領域リソースの生成・表示という中間的なプロセスを挿入することで、領域リソースを介したアノテーションの意味的グループ化をインターフェース上でも明示している。これは、5.4.2で論じた領域リソースとアノテーションの分離を、ユーザーの操作フローそのものに反映させる設計判断である。

特定のアノテーションが選択されると、アノテーション詳細パネルが表示される。詳細パネルは、5.4.4で導入したLinked Resourceの三分類に対応した「関連資料 Resources」（メディア）、「外部データ Linked Data」（典拠）、「参照文献 References」（書誌）の三つのタブに、地理座標を扱う「Location」タブを加えた計四つのタブから構成される（図~\ref{fig:ch5-annotation-panel}）。各タブでは、対応する種類のリソースを個別のダイアログを通じて追加・編集でき、追加の際にはRelation Hierarchyから関係タイプを選択する。この UIの構成は、Object全体に対するアノテーションと領域に対するアノテーションの双方で共通であり、詳細パネルの編集ハンドラは \texttt{test} コレクションの当該ドキュメントを直接更新することで、Object・Regionいずれの ID でも同一の操作を実現する。

\begin{figure}[htbp]
  \centering
  % \includegraphics[width=0.95\linewidth]{figures/chapter5/annotation-panel.png}
  \caption{アノテーションパネルと「関連資料 Resources」情報入力ダイアログ}
  \label{fig:ch5-annotation-panel}
\end{figure}

作成されたアノテーションは、RDF/Turtle形式およびIIIFマニフェスト形式の双方でエクスポートが可能である。この二形式並列出力は、5.4節で論じてきた設計原理を出力層に反映したものである。IIIF形式では、領域リソースは \texttt{oa:SpecificResource} として、その永続URIを \texttt{target.id} に含む形で出力される。同一 \texttt{regionId} を共有する複数のアノテーションは、それぞれが独立にvalidな \texttt{SpecificResource} として（\texttt{id}、\texttt{source}、\texttt{selector} を完全展開して）出力される。この冗長性は、既存のIIIFビューアとの互換性を担保するための設計判断であり、領域URIによる集約はRDF側に委ねられる。5.3において採用した「原表記の冗長保持」と同型の下位互換戦略である。Linked Resourceの詳細（書誌・典拠・メディアの内容や関係タイプ）は、IIIFには含めず、RDF側に委ねる。両形式は共通のURI体系を通じて相互参照される。

RDF形式では、\texttt{oa:Annotation} は \texttt{crmdig:D29\_Annotation\_Object} と、\texttt{oa:SpecificResource} は \texttt{crmdig:D35\_Area} と対応関係を持つ形で記述され、CIDOC CRMおよびCRMdigの語彙体系の上に位置づけられる。アノテーションの作成情報は \texttt{prov:wasAttributedTo}（作成主体）および \texttt{prov:generatedAtTime}（作成時点）によって、PROV-Oの語彙で記述される。各Linked Resourceの紐づけは、5.4.6で論じたようにE13 Attribute Assignment（\texttt{crm:E13\_Attribute\_Assignment}）として独立ノード化され、付与主体・時点・コメントを含む来歴情報を伴う。本文Markdownからは、5.4.7で論じたように \texttt{cito:cites} および \texttt{cito:discusses} 関係が自動抽出される。

このRDF出力を可視化した知識グラフの例を図~\ref{fig:ch5-knowledge-graph}に示す。この例では、IIIF準拠で提供される古代ローマ碑文3Dモデルの特定領域に付与された領域リソースを中心に、複数のプロジェクトからのアノテーションが集約されている構造が可視化される。領域リソースを軸として、アノテーションが複数のプロジェクト・ユーザーによって並立的に付与され、それぞれのアノテーションが Linked Resource として書誌・典拠・メディアを参照する構造が読み取れる。

\begin{figure}[htbp]
  \centering
  % \includegraphics[width=0.95\linewidth]{figures/chapter5/knowledge-graph.png}
  \caption{領域リソースを中心とした知識グラフの一例}
  \label{fig:ch5-knowledge-graph}
\end{figure}

\subsection{5.3との相補性と統合像}

本節で提示してきた史料側からの能動的アノテーションと、5.3で展開した二次文献からの抽出的アプローチは、最終的に同一の知識グラフ上で合流する。両者は排他的な選択肢ではなく、意味的接続を構成する二つの相補的な経路であり、両者を統合することによって、単独ではいずれも達成しえなかった精度と網羅性を備えた解釈ネットワークが実現する。

統合の要となるのは、参照される史料実体と、参照において用いられる書誌・典拠エンティティを、両アプローチが共通のURIで指し示す点である。5.3におけるInscription実体は、EDCS-IDを主キーとする外部データベースへの同定を経て安定した識別子を獲得していた。他方、本節におけるIIIFマニフェスト上のObjectは、それ自体としてマニフェストURIを持ち、さらに個々の物理的碑文としてはEDCS-IDあるいはWikidataのQIDといった同定子を介して外部と接続される。したがって、テキストを介して参照される対象としての碑文（5.3で扱った側面）と、物理的・視覚的対象としての碑文（本節で扱った側面）は、共通の識別子を介して同一の史料の二つの相として知識グラフ上で結び合わされる。同様に、書誌エンティティ（BibliographicResource）はCrossref DOIやOpenEdition書誌URIによって、典拠エンティティ（AuthorityResource）はWikidataのQIDやその他の権威データによって、両アプローチの間で共通の参照点を提供する。

この統合によって、たとえば以下のような横断的な問いに答えるための経路が開かれる。ある碑文の3Dモデル上の特定領域について、本節の枠組みで研究者が「この装飾要素はA氏の解釈によれば某地域の様式に属する」という書誌参照付きのアノテーションを付与したとする。同じ碑文に対して、5.3の抽出によって、Aとは異なるB氏の論文が「この碑文の年代はC世紀と考えられる」と論じている事実が既に構造化されているとすれば、当該碑文3DモデルのObjectノードから、二方向のアプローチによってもたらされた解釈の集積へと同時にアクセスできることになる。前者は5.4のRegion上のAnnotation経由で、後者は5.3のMention経由でグラフに登録されており、いずれも同じInscription URIに接続されているために、利用者はどちらの経路によって記述された情報かを意識することなく、当該史料に関する解釈の全体像を横断的に参照できる。

さらに重要なのは、この統合が単なる情報の並置にとどまらず、関係語彙の水準においても一貫している点である。5.3のRelation Assertionで用いた参照タイプ（内容説明・年代提示・翻刻文提供など）と、本節のCHAO v0.1におけるRelation Hierarchy（\texttt{mentions} / \texttt{depicts} / \texttt{contextualizes} 等）は、いずれも「史料に対して二次情報がどのような機能を果たすか」を記述する語彙であり、上位においては同じ意味論的空間に位置する。実際、5.3の参照タイプの多くは、CHAO v0.1のDirect Relationのサブクラスとして再解釈可能であり、両者を同一のオントロジー体系の中に統合する道筋は既に開かれている。したがって、二方向のアプローチによって蓄積される解釈は、単に同じ史料を指すだけでなく、同じ意味論的語彙のもとで比較・集約・検索可能な形で共存する。

また、この統合像はテキスト史料と非テキスト史料の境界を相対化する。5.4.1で述べたように、本節が採用した「部分領域への紐づけ」という設計は、非テキスト媒体において必然的に要請されるものであると同時に、テキスト史料の物理的側面や特定箇所への解釈にも同一の枠組みで対応しうる。実装上も、本研究のシステムはtextualエディタを通じてTEI構造化テキストの部分領域に対する同一枠組みでのアノテーション付与を可能にしている。これは、5.3で扱った碑文のテキスト（本文の翻刻）と、本節で扱った碑文の物理的対象（3Dモデル）とが、それぞれ異なる媒体でありながら同一の意味的接続基盤の上で扱われうることを示している。史料の媒体的多様性を横断する統一的な枠組みが、両方向のアプローチの統合を通じて具体化されるのである。

もっとも、この統合像はあくまで基盤の設計として提示されるものであり、実際に大規模なデータ蓄積のもとでどのような発見や解釈の展開が可能となるかは、今後の実証的研究に委ねられる。この点をはじめとする残された課題については、次章において改めて論じる。

\section{小括}

本章では、一次史料と二次文献の関係を、単純な参照リンクとしてではなく、意味と文脈を持つ学術的行為として表現する方法を検討した。既存のCIDOC CRMは対象と文書の基礎的接続を、CiTOは引用関係の型付けを、HuCIは人文学的引用における文脈保存の重要性を示していた。しかしながら、歴史研究における二次文献から一次史料への参照を扱うには、文献中の具体的な言及を独立した単位とし、領域固有の参照機能とその根拠を記述する必要がある。

そこで、Mentionを中心に、Textual Context、Corpus Reference、Inscription、Relation Assertionを接続するモデルを提示した。ローマ帝政アフリカ碑文学に関するOpenEdition上の四十の章・論文を対象とした実験では、2,796件の参照と1,685件の異なる碑文を抽出し、2,054件をEDCSへ同定した。参照の約三分の二が脚注に位置し、多くの言及が翻刻・間版間同定・年代・出土地といった複数の機能を同時に含んでいることは、文書構造と多値的な参照意味論を保持する必要性を示している。

さらに、この抽出的アプローチと相補的な位置づけとして、一次史料の側から二次情報を能動的に紐づける枠組みを検討した。この方向性は、参照する側の二次文献がテキストであることに依存する5.3の抽出的アプローチに対して、参照の起点を一次史料そのものに置くことで、非テキスト媒体をも含むより一般的な意味的接続を可能にするものである。「領域リソース」の導入によって注釈領域情報と解釈情報を分離し、複数主体による重層的な解釈の集約を可能にした。IIIF Presentation APIの拡張、CHAO v0.1におけるLinked Resourceの三分類とRelation Hierarchy、E13 Attribute Assignmentによる来歴管理、Markdown記法によるリソース内参照といった複数の設計要素を組み合わせ、これらをIIIF Semantic Editorとして実装した。

この二方向のアプローチによって、従来の「史料番号から文献を引く」索引は、「史料から言及箇所・文脈・参照機能・文献」へ、また逆方向に「史料の特定領域から書誌・典拠・メディアを介した多主体的解釈の蓄積」へと、双方向に拡張された。両者は共通のURI体系と関係語彙のもとで統合され、同一の知識グラフ上で相互に補完しあう。もっとも、デジタル化や参照形式、同定、意味分類の評価、大規模な運用における発見可能性の検証には課題が残る。これらの課題を、来歴情報と専門家検証を含む更新可能なワークフローとして扱うことが、今後の主要な課題である。
